\subsection{Optimization Procedure}
\label{app:procedure-algorithm}
The update loop below expands the main minimum-average-risk formulation, with candidate selection detailed in Appendix~\ref{app:procedure-mbr}. The retrieval corpus and task partitions remain fixed throughout optimization.
\begin{enumerate}
\item Initialize one predicate $p_{si}$ per skill, set $p_s\leftarrow p_{si}$, and evaluate the initial predicates on their retrieved validation tasks.
\item For iteration $i=1,\ldots,I$, save the iteration-start library $p^{(i)}$.
\item For each training task $q$ in fixed order, register it in the global seen-task history if it is new. Retrieve its candidates and evaluate their predicates. Materialize the bounded buckets for all retrieved or required catalog skills in sorted identifier order, then identify $\mathcal F_q$.
\item For each $s\in\mathcal F_q$, form the current-inclusive replay set, generate $M_p$ replacements, and evaluate all structurally valid replacements together with the unchanged predicate. Select an option with minimum average replay risk over this full set; the unchanged predicate ensures non-increase in average risk.
\item Commit the task's provisional repairs in deterministic skill order before processing the next task.
\item After the training pass, evaluate the candidate library on validation. For each changed skill, retain its candidate only if its validation error is strictly below the initial predicate's error; otherwise restore its iteration-start predicate.
\item Continue to the next iteration with this partially updated library. After iteration $I$, freeze the retained predicates for the final routing and downstream evaluations.
\end{enumerate}
